\documentclass[svgnames]{article}     % use "amsart" instead of "article" for AMSLaTeX format
%\geometry{landscape}                 % Activate for rotated page geometry

%\usepackage[parfill]{parskip}        % Activate to begin paragraphs with an empty line rather than an indent

\usepackage{graphicx}                 % Use pdf, png, jpg, or eps§ with pdflatex; use eps in DVI mode

%maths                                % TeX will automatically convert eps --> pdf in pdflatex
\usepackage{amssymb}
\usepackage{amsmath}
\usepackage{esint}
\usepackage{geometry}

% Inverting Color of PDF
\usepackage{xcolor}
\pagecolor[rgb]{0.19,0.19,0.19}
\color[rgb]{0.77,0.77,0.77}

%noindent
\setlength\parindent{0pt}

%pgfplots
\usepackage{pgfplots}

\usepackage{color}   %May be necessary if you want to color links
\usepackage{hyperref}
\hypersetup{
    colorlinks=true, %set true if you want colored links
    linktoc=all,     %set to all if you want both sections and subsections linked
    linkcolor=red,  %choose some color if you want links to stand out
}

%images
\graphicspath{{/Users/devaldeliwala/screenshots/}}                   % Activate to set a image directory

%tikz
\usepackage{pgfplots}
\pgfplotsset{compat=1.15}
\usepackage{comment}
\usetikzlibrary{arrows}
\usepackage[most]{tcolorbox}

\tcbset {
  base/.style={
    arc=0mm,
    bottomtitle=0.5mm,
    boxrule=0mm,
    colbacktitle=black!10!white,
    coltitle=black,
    fonttitle=\bfseries,
    left=2.5mm,
    leftrule=1mm,
    right=3.5mm,
    title={#1},
    toptitle=0.75mm,
  }
}

\definecolor{brandblue}{rgb}{0.34, 0.7, 1}
\newtcolorbox{mainbox}[1]{
  colframe=brandblue,
  base={#1}
}

\newtcolorbox{subbox}[1]{
  colframe=black!30!white,
  base={#1}
}

%Figures
\usepackage{float}
\usepackage{caption}
\usepackage{lipsum}


\title{Quantum Circuits}
\author{Deval Deliwala}
%\date{}                              % Activate to display a given date or no date

\begin{document}
\maketitle
%\section{}
%\subsection{}
\tableofcontents                     % Activate to display a table of contents
\newpage
This lesson introduces the \textit{quantum circuit} model of computation,
a standard description of quantum computations we'll use further onwards. \\

We will also introduce a few important mathematical concepts including
\textit{inner products}, the notion of \textit{orthogonality} and
\textit{orthonormality}, and \textit{projections} and \textit{projective
measurements}, which generalize standard basis measurements. Through these
concepts, we'll derive fundamental limitations on quantum information,
including the \textit{no-cloning theorem} and the impossibility to perfectly
discriminate non-orthogonal quantum states. \\

\section{Circuits}

In computer science, \textit{circuits} are models of computation in which
information is carried by wires via a network of \textit{gates}, which
represent operations that transform the information carried by the wires. \\

Although the word ``circuit" often refers to a circular path, circular paths
aren't actually allowed in most circuit models of computation. Quantum circuits
follow this pattern -- a quantum circuit represents a finite sequence of
operations that cannot contain feedback loops. 

\subsection{Boolean Circuits}

Here is an example of a classical Boolean circuit, where wires carry binary
values and the gates represent Boolean logic operations: 

\begin{figure}[H]
  \centering
    \includegraphics[width = 10cm]{screenshot.png}
    \caption*{XOR}
\end{figure}


The flow of information along the wires goes from left to right, the wires on
the left-hand side of the figure labeled $X$ and $Y$ are input bits, which can
each be set to whatever binary value we choose, and the wire on the right-hand
set is the output. The immediate wires take whatever values are determined by
the gates, which are evaluated from left to right. \\

The gates are AND gates ($\land$), OR gates ($\vee$), and NOT gates (¬). \\

The two small circles on the wires just to the right of the names $X$ and $Y$
represent \textit{fanout} operations, which simply create a copy of whatever
value is carried on the wire on which they appear, so that this value can be
input into multiple gates. \\

The particular circuit in this example computes the XOR, which is denoted by
the symbol $\oplus$

\begin{center}
\begin{tabular}{ l |r }
  $ab$ & $a \oplus b$ \\
   \hline
   00 & 0 \\ 
   01 & 1 \\ 
   10 & 1 \\ 
   11 & 0 \\ 
\end{tabular}
\end{center}

In the next diagram we consider just one choice for the inputs: $X = 0$ and $Y
= 1$. Each wire is labeled by value it carries so you can follow the
operations. The output value is $1$ in this case, which is the correct value
for $XOR: 0\oplus 1 = 1.$

\begin{figure}[H]
  \centering
    \includegraphics[width = 12cm]{screenshot 1.png}
\end{figure}

You can check the other three possible input settings in a similar way. 

\subsection{Other types of circuits}

In \textit{arithmetic circuits} the wires may carry integer values and the
gates may represent arithmetic operations, such as addition and multiplication.
We can also consider circuits that incorporate randomness, such as ones where
gates represent probabilistic operations. 

\subsection{Quantum Circuits}

In the quantum circuit model, wires represent qubits and gates represent
operations acting on these qubits. We'll focus for now on operations we've
encountered so far, namely \textit{unitary} operations and \textit{standard
basis measurements}. \\

Here is a simple example of a quantum circuit: 

\begin{figure}[H]
  \centering
    \includegraphics[width = 10cm]{screenshot 2.png}
\end{figure}

In this circuit we have a single qubit $X$ and a sequence representing unitary
operations on this qubit. Just like in the example above, the flow of
information goes from left to right. The first operation is Hadamard, then
\textit{S}, then another Hadamard, and a final $T$ operation. Applying the
entire circuit therefore applies the composition of these operations, $THSH$,
to the qubit $X$. \\ 

Sometimes we wish to explicitly indicate the input or output states to
a circuit. For example, if we apply the operation $TSHS$ to the state
$|0\rangle$, we obtain the state $\frac{1+i}{2}|0\rangle + \frac{1}{\sqrt{2}}|1\rangle $.
We may indicate this as follows: 

\begin{figure}[H]
  \centering
    \includegraphics[width = 12cm]{screenshot 3.png}
\end{figure}

Quantum circuits often have all qubits initialized to $|0\rangle $, but there
are also cases where we wish to set the input qubits to different states. \\

Here is how we can specify this circuit in Qiskit: 

\begin{figure}[H]
  \centering
    \includegraphics[width = 13cm]{screenshot 11.png}
\end{figure}


\begin{figure}[H]
  \centering
    \includegraphics[width = 13cm]{screenshot 12.png}
\end{figure}




Here is another example of a quantum circuit, this time with two qubits: 

\begin{figure}[H]
  \centering
    \includegraphics[width = 12cm]{screenshot 5.png}
\end{figure}

The second gate is a two-qubit gate -- a CNOT operation, where the solid
circle represents the control qubit and the circle with a plus inside denotes
the target qubit. \\

Before examining this circuit in greater detail and explaining what it does, it
is important to clarify how qubits are ordered in quantum circuits

\paragraph{Ordering of Qubits in Quantum Circuits} \mbox{} \\

The topmost qubit in a circuit has index 0, and corresponds to the rightmost
position in a Cartesian or Tensor Product. The second-to-top has index 1 and
corresponds to the position second-from-right in a Cartesian or Tensor Product,
and so on down to the bottom-most qubit which has the highest index, and
corresponds to the leftmost position in a Cartesian or Tensor product.\\

So for example, when we refer to a qubit in the zeroth position, we're referring
to the topmost qubit in a circuit diagram or the rightmost qubit in the
expressing of a quantum state vector; the qubit in the first position is
second-from-top in a circuit diagram or second from right in a quantum state
vector; and so on. This indexing convention is known as \textit{little-endian}
indexing, because the indexes start at the ``little end" when we think about
the significance of bits in binary representation of numbers.\\ 


Thus, in the circuit above, we are considering the circuit to be an operation
on two qubits $(X, Y)$. If the input is $|\psi\rangle |\phi\rangle $, then the
lower qubit $(X)$ starts in the state $|\psi\rangle $ and the upper qubit $(Y)$
starts in the state $|\phi\rangle .$\\

Now let's look at the circuit itself, moving from left to right through its
operations. \\

1. The first operation is a Hadamard Operation on $Y$. \\

When applying a gate to a single qubit like this, nothing happens to the other
qubits; nothing happening is equivalent to the identity operation. In our
circuit there is just one other qubit $X$, so the Hadamard Operation represents
this operation: 

\[
I \otimes H = \begin{pmatrix}
  \frac{1}{\sqrt{2}} & \frac{1}{\sqrt{2}} & 0 & 0 \\ 
  \frac{1}{\sqrt{2}} & -\frac{1}{\sqrt{2}} & 0 & 0 \\
  0 & 0 & \frac{1}{\sqrt{2}} & \frac{1}{\sqrt{2}} \\
  0 & 0 & \frac{1}{\sqrt{2}} & -\frac{1}{\sqrt{2}} 
\end{pmatrix} 
\] \vspace{5px}

2. The second operation is the controlled-NOT (CNOT) operation, where $Y$ is
   the control and $ X$ is the target. The CNOT gate's action on standard basis
   sates is as follows: 

\begin{figure}[H]
  \centering
    \includegraphics[width = 12cm]{screenshot 6.png}
\end{figure}

Given that we order the qubits as $(X, Y)$, the matrix representation of the
CNOT gate is this:

\[
\text{CNOT}  = \begin{pmatrix}
  1 & 0 & 0 & 0 \\ 0 & 0 & 0 & 1 \\ 0 & 0 & 1 & 0 \\ 0 & 1 & 0 & 0 
\end{pmatrix}
\] \vspace{5px}

The unitary operation of the entire circuit, which we'll call $U$ is the
composition of the two operations: 

\[
U = \begin{pmatrix}
  1 & 0 & 0 & 0 \\[4px] 0 & 0 & 0 & 1 \\[4px] 0 & 0 & 1 & 0 \\[4px] 0 & 1 & 0 & 0 
\end{pmatrix} \begin{pmatrix}
  \frac{1}{\sqrt{2}} & \frac{1}{\sqrt{2}} & 0 & 0 \\[4px] \frac{1}{\sqrt{2}} & -\frac{1}{\sqrt{2}} & 0 & 0 \\[4px] 0 & 0 & \frac{1}{\sqrt{2}} & \frac{1}{\sqrt{2}} \\[4px] 0 & 0 & \frac{1}{\sqrt{2}} & -\frac{1}{\sqrt{2}} 
\end{pmatrix} = \begin{pmatrix}
  \frac{1}{\sqrt{2}} & \frac{1}{\sqrt{2}} & 0 & 0 \\[4px] 0 & 0 & \frac{1}{\sqrt{2}} & -\frac{1}{\sqrt{2}} \\[4px] 0 & 0 & \frac{1}{\sqrt{2}} & \frac{1}{\sqrt{2}}  \\[4px] \frac{1}{\sqrt{2}} & -\frac{1}{\sqrt{2}} & 0 & 0
\end{pmatrix} 
\] \vspace{5px}

Recalling our knowledge of Bell states, we get that 

\begin{align*}
  &U|00\rangle = |\phi^+\rangle \\
  &U|01\rangle = |\phi^-\rangle \\
  &U|10\rangle = |\psi^+\rangle \\
  &U|11\rangle = -|\psi^-\rangle 
\end{align*}

This circuit gives us a way to convert the standard basis into the bell basis. (The -1 phase factor on the last state could be eliminated if we wanted by adding a controlled-Z gate at the beginning or a swap gate at the end, for instance).\\

In general quantum circuits can contain any number of qubit wires. We may also include classical bit wires which are indicated by double lines: 

\begin{figure}[H]
  \centering
    \includegraphics[width = 12cm]{screenshot 7.png}
\end{figure}

In this circuit we have a Hadamard Gate and a controlled-NOT gate on two qubits $X$ and $Y$, just like in the previous example. We also have two \textit{classical} bits, $A$ and $B$, as well as two measurement gates. The measurement gates represent standard basis measurements: the qubits are changed into their post measurement states, while the measurement outcomes are overwritten onto the classical bits to which the arrows point. \\

Here ins an implementation of this circuit using Qiskit: 

\begin{figure}[H]
  \centering
    \includegraphics[width = 12.5cm]{screenshot 9.png}
\end{figure}
\vspace{-25px}
\begin{figure}[H]
  \centering
    \includegraphics[width = 12.5cm]{screenshot 10.png}
\end{figure}


Sometimes it is convenient to depict a measurement as a gate that takes a qubit
as input and outputs a classical bit (as opposed to outputting the qubit its
post-measurement state and writing the result to a separate classical bit. This
means that the measured qubit has been discarded and can be safely ignored
thereafter. \\

For example, the following circuit diagram represents the same process as the
one in the previous diagram, but where we ignore $X$ and $Y$ after measuring
them. 


\begin{figure}[H]
  \centering
    \includegraphics[width = 12cm]{screenshot 12.png}
\end{figure}

\subsection{Diagrams of Qubit Gates}

Single-Qubit gates are generally shown as squares with a letter indicating
which operation it is, like this: 

\begin{figure}[H]
  \centering
    \includegraphics[width = 12cm]{screenshot 13.png}
\end{figure}

NOT (or $X$ ) gates are also sometimes denoted by a circle around a + sign. 

Swap gates are denoted as follows: 

\begin{figure}[H]
  \centering
    \includegraphics[width = 12cm]{screenshot 14.png}
\end{figure}

Controlled-gates, meaning that gates that describe controlled-unitary
operations, are denoted by a filled-in circle (indicating the control)
connected by a vertical line to whatever operation is being controlled.  For
instance, controlled-NOT gates, controlled-controlled-NOT gates, and
controlled-swap gates are denoted like this: 

\begin{figure}[H]
  \centering
    \includegraphics[width = 12cm]{screenshot 15.png}
\end{figure}

Arbitrary unitary operations on multiple qubits may be viewed as gates. They
are depicted by rectangles labeled by the name of the unitary operation. For
instance, here is a depiction of an (unspecified) unitary operation $U$ as
a gate, along with a controlled version of this gate: 

\begin{figure}[H]
  \centering
    \includegraphics[width = 12cm]{screenshot 16.png}
\end{figure}

\section{Inner Products, Orthonormality, and Projections}


The notions of the \textit{inner product, }\textit{orthogonality}, and \textit{orthonormality} for sets of
vectors, and \textit{projection} matrices will allow us to introduce a handy
generalization of standard basis measurements. 


\subsection{Inner Products}

We refer to an arbitrary column vector as a ket, such as 

\[
|\psi\rangle = \begin{pmatrix}
  \alpha_1 \\ \alpha_2 \\ \vdots \\ \alpha_n
\end{pmatrix} 
\] \vspace{5px}

the corresponding bra vector is the \textit{conjugate transpose} of this
vector: 

\[
\langle \psi | = (|\psi\rangle)^\dagger = (\bar\alpha_1 \quad \bar\alpha_2
\quad \cdots \quad \bar\alpha_n)
\] \vspace{5px}

Alternatively if we have some classical state set $\Sigma$ in mind, and we
express a column vector as a ket, such as 

\[
  |\psi\rangle = \sum_{a\in\Sigma} \alpha_a |a\rangle 
\] \vspace{5px}

then the corresponding bra vector is the conjugate transpose

\[
\langle \psi | = \sum_{a\in\Sigma} \bar\alpha_a \langle a | 
\] \vspace{5px}

We also observed that the product of a bra and ket vector, viewed as matrices
either having a single row or a single column, results in a scalar.
Specifically, if we have two (column) vectors

\[
|\psi\rangle = \begin{pmatrix}
  \alpha_1 \\ \alpha_2 \\ \vdots \\ \alpha_n
\end{pmatrix} \quad \text{and} \quad |\phi\rangle = \begin{pmatrix}
  \beta_1 \\ \beta_2 \\ \vdots \\ \beta_n
\end{pmatrix} 
\] \vspace{5px}
then, 

\[
\langle \psi | \phi \rangle = \langle \psi | |\phi\rangle = (\bar\alpha_1 \quad \bar\alpha_2 \quad \cdots \quad \bar\alpha_n) \begin{pmatrix}
  \beta_1 \\ \beta_2 \\ \vdots \\ \beta_n
\end{pmatrix} = \bar\alpha_1 \beta_1 + \cdots + \bar\alpha_n \beta_n
\] \vspace{5px}

Alternatively, if we have two column vectors that we have written as 

\[
|\psi\rangle = \sum_{a\in\Sigma} \alpha_a |a\rangle \quad \text{and} \quad |\phi\rangle = \sum_{b\in\Sigma} \beta_b |b\rangle 
\] \vspace{5px}

then, 

\begin{mainbox}{Inner Product Definition}
\begin{align*}
  \langle \psi | \phi \rangle &= \langle \psi | |\psi\rangle \\
                              &= \left( \sum_{a\in\Sigma} \bar\alpha_a \langle a |\right) \left( \sum_{b\in\Sigma} \beta_b |b\rangle  \right) \\ 
                              &= \sum_{a\in\Sigma} \sum_{b\in\Sigma} \bar\alpha_a \beta_b \langle a | b\rangle \\ 
                              &= \sum_{a\in\Sigma} \bar\alpha_a \beta_a
\end{align*}
\end{mainbox}

where the last equality stems from the observation that $\langle a | a \rangle = 1$ and $\langle a | b \rangle = 0 $ for classical states $a$ and $b$ satisfying $a \neq b$. \\ 

The value $\langle \psi | \phi \rangle $ is called the \textit{inner product} between vectors $|\psi\rangle $ and $|\phi\rangle $. Inner products are critically important in quantum information and computation. \\

Let us go over some basic facts about the inner product. 

\paragraph{Relationship to the Euclidean Norm} \mbox{} \\

The inner product of any vector with itself, 

\[
\langle \psi | \psi \rangle = || |\psi\rangle ||^2. 
\] \vspace{5px}

Thus, the euclidean norm of a vector may be expressed as 

\[
|||\psi\rangle || = \sqrt{\langle \psi | \psi \rangle }
\] \vspace{5px}

And since the euclidean norm can only equal 0 if every entry is equal to 0, for every vector $|psi\rangle $ we have 

\[
\langle \psi | \psi \rangle \geq 0
\] \vspace{5px}

\paragraph{Conjugate Symmetry} \mbox{} \\ 

For any two vectors 

\[
|\psi\rangle  = \sum_{a\in\Sigma} \alpha_a |a\rangle \quad \text{and} \quad |\phi\rangle = \sum_{b\in\Sigma} \beta_b |b\rangle 
\] \vspace{5px}

we have 

\[
\langle \psi | \phi \rangle = \sum_{a\in\Sigma} \bar\alpha_a \beta_a \quad \text{and} \quad \langle \phi | \psi \rangle = \sum_{a\in\Sigma} \bar\beta_a\alpha_a
\] \vspace{5px}

and therefore 


\[
  \overline{\langle \psi | \phi \rangle} = \langle \phi | \psi \rangle 
\] \vspace{5px}

\paragraph{Linearity in the Second Argument (and conjugate linearity in the first)} \mbox{} \\

Let us suppose that $|\psi\rangle , |\phi_1\rangle , and |\phi_2\rangle $ are vectors and $\alpha_1$ and $\alpha_2$, are complex numbers. If we define a new vector 

\[
|\phi\rangle = \alpha_1|\phi_1\rangle  + \alpha_2|\phi_2\rangle, 
\] \vspace{5px}

then 

\[
\langle \psi | \phi \rangle = \langle \psi | (\alpha_1|\phi_1\rangle + \alpha_2 |\phi_2\rangle ) = \alpha_1\langle \psi | \phi_1 \rangle + \alpha_2\langle \psi | \phi_2 \rangle 
\] \vspace{5px}


The inner product is \textit{linear} in the second argument. This can be verified either through the formulas above or simply by noting that matrix multiplication is \textit{linear} in each argument. \\

Combining this fact with conjugate symmetry reveals that the inner product is \textit{conjugate linear} in the first argument. That is, if $|\psi_1\rangle, |\psi_2\rangle, \text{ and } |\phi\rangle $ are vectors and $\alpha_1$ and $\alpha_2$ are complex numbers, and we define

\[
|\psi\rangle = \alpha_1 |\psi_1\rangle + \alpha_2 |\psi_2\rangle 
\] \vspace{5px}

then 

\[
\langle \psi | \phi \rangle  = (\bar\alpha_1 \langle \psi_1 | + \bar\alpha_2 \langle \psi_2 | )|\phi\rangle = \bar\alpha_1 \langle \psi_1 | \phi \rangle + \bar\alpha_2 \langle \psi_2 | \phi \rangle  
\] \vspace{5px}

\paragraph{The Cauchy-Schwarz Inequality} \mbox{} \\

For every choice of vectors $|\phi\rangle $ and $|\psi\rangle $ having the same
number of entries, we have 

\[
  | \langle \psi | \phi \rangle | \leq \lVert \psi\rangle \rVert \lVert |\phi\rangle \rVert
\] \vspace{5px}

This inequality is super important and useful. 

\subsection{Orthogonal and Orthonormal Sets}

Two vectors $|\phi\rangle $ and $|\psi\rangle $ are said to be
\textit{orthogonal} if their inner product is 0: 

\[
\langle \psi | \phi \rangle = 0
\] \vspace{5px}

A set of vectors $\{ |\psi_1\rangle, ... , |\psi_m\rangle \} $ is an
\textit{orthogonal set} if every vector in the set is orthogonal to every other
vector in the set: 

\[
\langle \psi_j| \psi_k \rangle = 0
\] \vspace{5px}

for all choices of $j, k \in \{1, ... , m\}$ for which $j \neq k$. \\ 

A set of vectors $\{ |\psi_1\rangle , ... , |\psi_m\rangle \} $ is called an
\textit{orthonormal set} if it is an orthogonal set and, in addition every vector is a unit vector: 

\[
  \langle \psi_j | \psi_k \rangle = 
  \begin{cases}
    1 \quad j = k \\ 0 \quad j \neq k
  \end{cases}
\] \vspace{5px}

Finally, a set $\{ |\psi_1\rangle, ... , |\psi_m\rangle\} $ is an
  \textit{orthonormal basis} if, in addition to being an orthonormal set, it
  forms a basis. This is equivalent to being both an orthonormal set and $m$ 
  being equal to the dimension of the space from which the vectors are drawn.\\

For example, for any classical state set $\Sigma$, the set of all standard
basis vectors

\[
  \{ |a\rangle : a\in\Sigma \}
\] \vspace{5px}

is an orthonormal basis. The set $\{|+\rangle, |-\rangle \}$ is an orthonormal
basis for the 2-dimensional space corresponding to a single qubit, and the Bell
Basis $\{|\phi^+\rangle , |\phi^-\rangle , |\psi^+\rangle , |\psi^-\rangle $ is
  an orthonormal basis for the 4-dimensional space corresponding to two qubits.






\end{document}
